\chapter{Notes on Isaiah}\label{app:isaiah}
Most of the prophetic writings are either in judgment or hope for the people; they are not usually neutral. One theory postulates that there are two Isaiahs: one comprising Isaiah~1--39 and dealing with judgment, the other Isaiah~40--66 and dealing with hope (and possibly a third comprising just 56--66). However, there is plenty of hope in the first section and plenty of judgment in the second section which may controvert this theory.

\section{Isaiah 2:16}\label{app:isa-2-16}
Sailing in antiquity was a dangerous activity~--- it was taking your life in your hands. Sailors would go as close to the coast as possible so that if the ship breaks up (cf. Jonah 1:4) they could live through the incident by swimming to shore.

The \emph{Hebrew and Aramaic Lexicon of the Old Testament} by Koehler and Baumgartner defines \Hebrew{שְׂכִיָּה} as follows:
\begin{quote}
    Ug. \emph{$\underline{t}$kt} (Gordon \emph{Textbook} \S19:2680; Aistleitner 2862; Driver \emph{Myths}$^2$ 160a; cf. Fisher \emph{Parallels} 2: p. 8 entry 5) < Eg. \emph{\'skty} ship (Erman-G. 4:315), see Lambdin \emph{Loan Words} 154f; Ellenbogen \emph{Foreign Words} 154; cf. also Wildberger BK 10:94: \textbf{ship} Is 2$_{16}$. $\dagger$
\end{quote}

The following is from Ellenbogen's \emph{Foreign Words in the Old Testament}:
\begin{quotation}
    {\noindent\Hebrew{שכיות} -- Ships~~~~~~~~~~~~~~~~~~~~~~~~~~~Egyptian -- $\underline{\emph{\'sk.tj}}$ \includegraphics[scale=1.1]{images/egt-ship} (ship)}
    
    {\noindent Isa. 2:16 -}
    
    The LXX translation of the phrase \Hebrew{החמדה כל־שכיות ועל}, which runs: \Greek{ἐπὶ πᾶσαν θέαν πλοίων κάλλους}, is very remarkable; it incorporates both what seems to be the real meaning of the term \Hebrew{שכיות}, namely, \Greek{πλοιον} ``ship,'' and also what is required by the traditional understanding of the word, namely \Greek{θέᾶ} ``sight, view.'' The Vulgate renders \Hebrew{שכיות} by $\underline{\text{visus}}$,\footnote{Prob. ``views''} and the Peshi\d{t}ta by \Syriac{ܕܘܩܐ} ``view.'' Rashi has a note saying that \Hebrew{שכיות} refers to palaces whose floors are paved with marble mosaics. This seems to be purely conjectural, and possibly based on the Targum's rendering of the term (\Hebrew{בירנתא} ``castle'').
    
    The Verses Isa. 2:13, 14, 15, and 17 contain parallelisms, so it would appear reasonable to assume that Verse 2:16 also contains a parallelism that would match the word \Hebrew{אניות} ``ships,'' so Budde-Begrich ($\underline{\text{ZATW}}$ 49, p.198) are apparently right in seeing in Hebrew \Hebrew{שכיות} the borrowed Egyptian term $\underline{\emph{\'sk.tj}}$, a word which is attested from the days of the New Kingdom. [See Erman-Grapow, $\underline{\text{WB}}$ IV, p.315.] The Egyptian derivation of \Hebrew{שכיות} was also accepted by Albright ($\underline{\text{Bertholet Festschrift}}$, p.5) who mentions that H.L. Ginsberg identifies Ugaritic $\underline{\text{\underline{t}kt}}$ with Hebrew \Hebrew{שכיות}.
\end{quotation}

Note that the Septuagint allows for the rendering found in the King James Version.
